\begin{table*}[t]
\centering
\small
\resizebox{\textwidth}{!}{
\begin{tabular}{l l c l l c}
\toprule
\multirow{2}{*}{\textbf{Model}} & \textbf{Architecture} & \multirow{2}{*}{\textbf{Loss}} & \textbf{SSL Pre-training} & \textbf{Phone Recognition} & \multirow{2}{*}{\textbf{Langs}} \\
 & \textbf{(Enc / Dec)} & & \textbf{Data} & \textbf{Data} & \\
\midrule

\lvsixty \cite{xu22b_interspeech} & Enc: Wav2Vec2 & CTC & LibriLight & MLS, CV, Babel & 40+ \\
\xlsrf \cite{xu22b_interspeech} & Enc: XLSR53 & CTC & MLS, CV, Babel & MLS, CV, Babel & 40+ \\
\mipa \cite{taguchi23_interspeech} & Enc: XLSR53 & CTC & MLS, CV, Babel & CV 11.0 & 7 \\

\midrule

\zipactc \cite{zipa} & Enc: Zipformer & CR-CTC & None & IPAPack++ & 88 \\
\zipactcns \cite{zipa} & Enc: Zipformer & CR-CTC & None & IPAPack++ \& PL & 4k \\

\midrule

\multirow{2}{*}{\powsm \cite{powsm}} & Enc: E-Branchformer & \multirow{2}{*}{CTC-Att} & \multirow{2}{*}{None} & \multirow{2}{*}{IPAPack++} & \multirow{2}{*}{88} \\
 & Dec: Transformer & & & & \\

\powsmctc (\href{https://huggingface.co/espnet/powsm_ctc}{ours}) & Enc: E-Branchformer & Int-CTC & None & IPAPack++ & 88 \\

\midrule

\gemini \cite{comanici2025gemini} & Closed & N/A & Closed & Closed & >200 \\

\multirow{2}{*}{\qweni \cite{xu2025qwen3omnitechnicalreport}} & Enc: AuT & \multirow{2}{*}{AR} & \multirow{2}{*}{Closed} & \multirow{2}{*}{Closed} & \multirow{2}{*}{19} \\
 & Dec: MoE Transformer & & & & \\


\bottomrule
\end{tabular}
}
\caption{Included PR systems. 
Architecture abbrv.: Encoder (Enc), Decoder (Dec), Audio Transformer (AuT), Mixture-of-Experts (MoE);
Loss abbrv.: Consistency Regularized CTC (CR-CTC) \cite{crctc}, Hybrid CTC/Attention (CTC-Att) \cite{watanabe2017hybridctc}, Intermediate CTC (Int-CTC) \cite{lee2021intermediate}, Autoregressive (AR);
Data abbrv.: Multilingual LibriSpeech (MLS), Common Voice (CV), Pseudo-labeled (PL).}
\label{tab:model-summary}
\end{table*}


\section{Evaluation Framework of \bench}
\bench covers intrinsic (\autoref{ss:core}) and extrinsic (\autoref{ss:downstream}) evaluations shown in \autoref{fig:phonebench-design}. 
Intrinsic evaluation compares predicted transcriptions to gold labels, while extrinsic evaluation measures transcriptions and internal representations on downstream tasks.
In extrinsic evaluation, transcriptions provide a direct and interpretable signal of explicit phonetic content, whereas representations are commonly used in downstream tasks but may encode non-phonetic information.
\autoref{tab:evaltasks} summarizes included datasets and metrics.


\subsection{Intrinsic: Core Capability}\label{ss:core}
We use \textbf{Phonetic Feature Error Rate (PFER)} to measure the distance between reference and predicted transcriptions. 
Unlike Phone Error Rate (PER), which treats each phone as a token, PFER computes the edit distance $D(\cdot, \cdot)$ over articulatory features $\text{feat}(\cdot)$ such as roundness or voicing. As shown in \autoref{eq:pfer}, where $u$ denotes an utterance (sequence of phones where $i$ indexes this sequence) and $u^*$ its ground truth, PFER is calculated as the total feature edit distance across all utterances divided by the total number of phones, representing the percentage of incorrect features \cite{panphon}. 
\begin{equation}
\label{eq:pfer}\small
    \text{PFER} = \frac{1}{\sum_i |u_i^*|}\sum_i D(\text{feat}(u^*_i),\text{feat}(u_i))
\end{equation}

The tasks comprise two categories: 
\textbf{Variation of seen languages} includes regional and non-native speech, testing whether PR systems rely excessively on seen patterns rather than the actual input.
\textbf{Unseen languages} assess the system's language-agnostic phonetic knowledge. 
Details of each task and dataset are in \autoref{appendix:datadetails-pr}.


\subsection{Extrinsic: Downstream Utility}\label{ss:downstream}
We evaluate PR systems using two downstream probes, namely the transcript probe and the representation probe.
For \textbf{transcript probe (TP)}, input consists of predicted phonetic transcriptions, and the probe is a text-based bi-GRU.
For \textbf{representation probe (RP)}, following the setup in \citet{hear}, we use the last layer's hidden representations as input and temporal pooling with attention followed by a Multi-Layer Perceptron as the probe.
Metrics for each task are listed in \autoref{tab:evaltasks} and the detailed experimental setup is in \autoref{appendix:setup}.

We consider three categories of downstream tasks where phonetic information is essential. 
In \textbf{pathological speech assessment}, phonetic transcriptions are used to document patients' speech and support diagnosis and treatment planning \cite{ball2009importance, nelson2020use}. 
In \textbf{L2 speech assessment}, phonetic cues enable pronunciation feedback \cite{Franco2010EduSpeakAS} and accent classification \cite{angkititrakul2006phone-based-accent}. 
In \textbf{multilingual speech identification}, analyzing phonetic and phonological differences across languages and dialects, such as phone inventories, phonotactics, and phoneme realization, is crucial \cite{schultz2006multilingualsp}.
We describe each task and dataset in detail in \autoref{appendix:datadetails-downstream}.
